% answers_short.tex —— 2007 年·普通高等学校招生考试·新课标卷·文科 · 简短答案速查
% 由 tools/make_short_answers.py 自 answers.tex 的 \gksolhead/\gkanswer 提取，
% 并与 2007/2007_数学真题_新课标文_答案要点.md 逐题比对后生成，勿手改（改源 MD 后重跑）。

\begin{gknotice}
\textbf{使用说明}：本节只给最终答案便于快速核对，完整推导见解析版。答案为编写者独立推导，\textbf{非官方标准答案}。
\end{gknotice}

\gkpart{选择题}
\begin{center}\begin{tabular}{ccccccccccccc}
\toprule
题号 & 1 & 2 & 3 & 4 & 5 & 6 & 7 & 8 & 9 & 10 & 11 & 12 \\
\midrule
答案 & A & C & A & D & C & B & C & B & C & D & D & B \\
\bottomrule
\end{tabular}\end{center}

\gkpart{填空题}
\begin{itemize}[leftmargin=2.6em]
  \item[\textbf{13.}] $3$。
  \item[\textbf{14.}] $-1$。
  \item[\textbf{15.}] $4-4\mathrm i$。
  \item[\textbf{16.}] $\dfrac12$。
\end{itemize}

\gkpart{解答题}
\begin{itemize}[leftmargin=2.6em]
  \item[\textbf{17.}] $AB=\dfrac{s\sin\beta\tan\theta}{\sin(\alpha+\beta)}$。
  \item[\textbf{18.}] （1）$CD=2$；（2）总有 $AB\perp CD$。
  \item[\textbf{19.}] （1）增区间 $\left(-\dfrac32,-1\right)$，$\left(-\dfrac12,+\infty\right)$，减区间 $\left(-1,-\dfrac12\right)$；（2）最大值 $\ln\dfrac72+\dfrac1{16}$，最小值 $\ln2+\dfrac14$。
  \item[\textbf{20.}] （1）$\dfrac34$；（2）$\dfrac23$。
  \item[\textbf{21.}] （1）$-\dfrac34<k<0$；（2）不存在。
  \item[\textbf{22.}] （1）证明见解析；（2）$90^\circ$。
\end{itemize}

\gkpart{选做题}
\begin{itemize}[leftmargin=2.6em]
  \item[\textbf{23.}] （1）$(x-2)^2+y^2=4$，$x^2+(y+2)^2=4$；（2）$x+y=0$。
  \item[\textbf{24.}] （1）$x<-7$ 或 $x>\dfrac53$；（2）$-\dfrac92$。
\end{itemize}
